\section{Introduction}
\label{sec:intro}

A BPE tokenizer is a list of merges~\citep{gage1994new,sennrich2016neural}.
Training adds merges to this list one at a time, each time choosing the most frequent adjacent pair.
At encoding time, the list order becomes a priority: the tokenizer repeatedly applies the highest-priority merge that fits~\citep{radford2019language}.
The tokenizer is then frozen into every model built on top of it.
Because the frequency rule is so standard, merge order is almost never treated as a design choice.
Our main question is simple: {\bf does the order of BPE merges change the quality of a model built on the tokenizer?}

The question matters in practice.
Practitioners change merge order whenever they re-weight a tokenizer corpus toward a domain, add domain-specific tokens, or merge vocabularies from different sources.
If quality is sensitive to merge order, these changes need care.
If models are robust to it, merge order can be changed freely, for example for domain adaptation.
Prior work compares tokenizer \emph{algorithms}~\citep{bostrom2020byte,ali2024tokenizer}, vocabulary sizes and compression~\citep{galle2019investigating,goldman2024unpacking}, intrinsic efficiency measures~\citep{zouhar2023tokenization}, stochastic segmentation~\citep{provilkov2020bpe,kudo2018subword}, and domain vocabularies~\citep{beltagy2019scibert,gu2021domain}.
However, these comparisons change many things at once.
We are not aware of a controlled study that varies \emph{only} merge order while holding data, vocabulary size, model, and training text fixed.
Prior work also does not separate \emph{which} merges are learned from the \emph{priority} with which they are applied.

We address this gap with a controlled from-scratch study (\figref{fig:overview}).
We train six byte-level BPE tokenizers with 8,192 vocabulary entries on the same corpus, and we vary merge order in two ways.
Four tokenizers change the \emph{learning order}, and hence the vocabulary:
standard frequency BPE (\freq), domain-weighted frequency (\domain), random choice among frequent pairs (\randlearn), and least-frequent-first (\revlearn).
Two tokenizers change only the \emph{application order}: they keep \freq's exact merges and vocabulary, but apply them in a random (\randorder) or reversed (\revorder) priority.
For each tokenizer, we pre-train identical 23M-parameter GPT-style models on the same 200\,MB of Wikipedia and PubMed text.
We measure tokenizer-independent bits-per-byte (\bpb), and then fine-tune each model on \sst, \agnews, and \pubmedrct sentence classification.

The results split cleanly by metric.
For language modeling, merge order clearly matters.
On general text, \freq beats every alternative.
The four non-frequency orders raise \bpb by 1.8--5.8\% on general text and by 2.8--9.8\% on biomedical text, and all 10 \bpb comparisons against \freq are significant after Holm correction.
The same-vocabulary controls are the most telling.
\randorder and \revorder lose 1.8--5.7\% \bpb with \emph{exactly} the same merges as \freq, because short merges that fire early consume their inputs and block long, frequent tokens.
For example, \revorder splits ``\,hypertension'' into eight pieces, while \freq keeps it as one token.
Compression explains the whole ranking: bytes-per-token rank-predicts general-text \bpb perfectly (Spearman $\rho=-1.0$, $n=6$).
For fine-tuned classification, in contrast, every tokenizer lands within $-0.8\%$ to $+1.6\%$ of \freq, and none of the 15 comparisons is significant.
Even a nearly character-level tokenizer (\revlearn) is within 1\% of \freq on every task.
Domain weighting buys a real but small domain gain: 0.8\% \bpb on PubMed, at a cost of 0.4\% on general text.

\begin{figure}[t]
    \centering
    \resizebox{0.95\linewidth}{!}{%
    \begin{tikzpicture}[
        box/.style={draw, rounded corners=2pt, align=center, font=\small, minimum height=0.9cm, inner sep=4pt},
        tok/.style={box, fill=blue!6, text width=2.25cm},
        fam/.style={box, fill=gray!12, font=\small\bfseries, text width=3.0cm},
        outbox/.style={box, fill=green!8, text width=3.2cm},
        arr/.style={-{Stealth[length=2mm]}, thick}
    ]
    \node[box, fill=yellow!15, text width=2.6cm] (corpus) at (0,0) {Tokenizer corpus\\{\footnotesize 50\,MB Wiki + PubMed}};
    \node[fam] (learn) at (3.4,1.1) {Learning order\\{\footnotesize\normalfont which merges are learned}};
    \node[fam] (apply) at (3.4,-1.1) {Application order\\{\footnotesize\normalfont same merges as \freq, new priority}};
    \node[tok] (t1) at (6.6,1.65) {\freq, \domain};
    \node[tok] (t2) at (6.6,0.55) {\randlearn, \revlearn};
    \node[tok] (t3) at (6.6,-1.1) {\randorder, \revorder};
    \node[box, fill=orange!10, text width=2.7cm] (lm) at (9.9,0) {Identical 23M LM\\{\footnotesize same 200\,MB text, 2 seeds}};
    \node[outbox] (bpb) at (13.4,0.75) {Held-out bits-per-byte\\{\footnotesize general / biomedical}};
    \node[outbox] (ft) at (13.4,-0.75) {Fine-tuned accuracy\\{\footnotesize \sst, \agnews, \pubmedrct}};
    \draw[arr] (corpus.east) -- ++(0.2,0) |- (learn.west);
    \draw[arr] (corpus.east) -- ++(0.2,0) |- (apply.west);
    \draw[arr] (learn.east) -- (t1.west);
    \draw[arr] (learn.east) -- (t2.west);
    \draw[arr] (apply.east) -- (t3.west);
    \draw[arr] (t1.east) -- (lm.west);
    \draw[arr] (t2.east) -- (lm.west);
    \draw[arr] (t3.east) -- (lm.west);
    \draw[arr] (lm.east) -- (bpb.west);
    \draw[arr] (lm.east) -- (ft.west);
    \end{tikzpicture}}
    \caption{Study design. We train six byte-level BPE tokenizers (8,192 entries each) that differ only in merge order. Four change the \emph{learning order}, which changes the vocabulary. Two keep \freq's exact merges and change only the \emph{application order}, \ie merge priority at encoding time. Every tokenizer is paired with an identical language model trained on the same text, and we evaluate it with tokenizer-independent bits-per-byte and with fine-tuned classification.}
    \label{fig:overview}
\end{figure}

In summary, we make the following contributions:
\begin{itemize}[leftmargin=*,itemsep=0pt,topsep=0pt]
    \item We design a controlled study that isolates BPE merge order, and we separate learning order (which merges exist) from application order (which merges fire first) through same-vocabulary controls.
    \item We show that merge order significantly affects language modeling (+1.8--9.8\% \bpb) and that pure reordering of an identical vocabulary costs up to 5.7\%, leaving 59--79\% of the vocabulary unused.
    \item We identify compression as the mechanism: bytes-per-token rank-predicts \bpb perfectly on general text ($\rho=-1.0$) and strongly on biomedical text ($\rho=-0.94$), while R\'enyi efficiency does not.
    \item We find that fine-tuned classification is robust to merge order: our confidence intervals rule out effects larger than about 1.7\% on \agnews and \pubmedrct, and about 3.5\% on \sst.
\end{itemize}

The rest of the paper is organized as follows.
\Secref{sec:related} reviews related work.
\Secref{sec:method} describes the tokenizers, models, and evaluation.
\Secref{sec:results} presents our results, and \secref{sec:discussion} discusses their implications and limitations.
